\documentclass[10pt]{article}
\usepackage[letterpaper,margin=0.66in]{geometry}
\usepackage{microtype}
\usepackage{amsmath,amssymb}
\usepackage{booktabs}
\usepackage{enumitem}
\usepackage{xcolor}
\usepackage{hyperref}
\hypersetup{colorlinks=true,linkcolor=blue!45!black,citecolor=blue!45!black,urlcolor=blue!45!black,pdftitle={The D-Score Snapshot Alignment Theorem},pdfauthor={Justin D. Hart}}
\setlength{\parindent}{0pt}
\setlength{\parskip}{2.5pt}
\setlist{nosep,leftmargin=*}
\raggedbottom
\emergencystretch=2em

\newcommand{\Var}{\operatorname{Var}}
\newcommand{\E}{\mathbb{E}}

\title{\vspace{-1.2em}\textbf{The D-Score Snapshot Alignment Theorem}\\
\large A delay-variance budget for asynchronous biodiversity indicators}
\author{Justin D. Hart\\\small Viridis Research}
\date{5 August 2026\\\small Final reconciled manuscript}

\begin{document}
\maketitle

\begin{abstract}
Biodiversity composites often fuse observations collected on incompatible
clocks: satellite products may refresh daily, field inventories seasonally,
and genetic surveys over years.  A weighted average can therefore appear
precise while its components describe different ecological states.  We isolate
this temporal-coherence problem in a local common-mode model.  For normalized
nonnegative effective weights $w_i$, observation delays $\tau_i$, and harmonic
frequency $\omega$, the fused response has alignment efficiency
\(
\eta(\omega)=\left|\sum_i w_i e^{-\mathrm{i}\omega\tau_i}\right|^2.
\)
We derive the exact pairwise loss identity and the conservative bound
\(
\eta(\omega)\ge 1-\omega^2\Var_w(\tau).
\)
Consequently, a declared ecological bandwidth $\Omega$ and target coherence
$q$ imply an auditable sufficient rule,
\(
\Var_w(\tau)\le(1-q)/\Omega^2.
\)
Equal delays preserve perfect coherence, whereas two equally weighted channels
separated by half a cycle cancel completely.  A preregistered standard-library
Python gate passed 21 of 21 analytic, randomized, convergence, negative-control,
and time-domain checks.  Separately, Aristotle returned six Lean theorems that
an independent local audit rebuilt in Lean 4.28.0/Mathlib, with no proof holes,
forbidden constructs, or axioms beyond \texttt{propext},
\texttt{Classical.choice}, and \texttt{Quot.sound}.  The result is a formally
verified conditional signal-model statement, not an empirical validation of
D-Score, a novelty or priority finding, or evidence of field utility.
\end{abstract}

\textbf{Keywords:} biodiversity monitoring; data fusion; temporal alignment;
composite indicators; D-Score; sensor synchronization; measurement coherence

\section{Introduction}

Composite biodiversity indicators are assembled from measurements that differ
in biological meaning, spatial support, error structure, sampling protocol, and
time.  The last difference is easy to hide.  A species inventory dated 2021, a
satellite productivity layer dated 2026, and a genetic-diversity estimate dated
2018 can be placed in the same row of a database and averaged without any
arithmetic warning.  The result is numerically well formed but may not describe
an ecosystem that ever existed at one time.

The D-Score research program makes this concern concrete.  Its validation
protocol combines taxonomic, phylogenetic, genetic, and environmental inputs.
Prior Viridis work has separately studied component priorities, finite-sample
measurement floors, spatial allocation of monitoring power, optimal revisit
cadence, and weights robust to strategic gaming.  None of those results answers
the narrower question here: \emph{when do component timestamps permit a
weighted score to be interpreted as a coherent snapshot?}

The broader literature supplies both motivation and caution.  Multi-sensor
fusion systems require a shared time reference and explicit temporal sample
alignment; experimental work shows that software synchronization can prevent
cumulative desynchronization across heterogeneous sensors
\cite{piatkowski2025}.  Biodiversity data products similarly combine remote
sensing, field surveys, and other observation streams, for which fusion can add
substantial information \cite{schulte2018}.  Yet ecological timing is not merely
an instrumentation nuisance.  Global phenology maps reveal spatially displaced
seasonal cycles and validate them against independent products
\cite{terasakihart2025}; nearly half a million plant and animal time series show
divergent phenological shifts \cite{lang2025}.  Real biological asynchrony can be
signal rather than error.  Monitoring uncertainty also limits the detectable
evidence for biodiversity trend change even at large sample sizes
\cite{leung2024}.

We therefore seek a deliberately bounded result.  We do not claim that all
D-Score components share one dynamic, that all delays should be removed, or that
temporal coherence proves ecological validity.  We ask what follows \emph{if}
a set of normalized channels is intended to measure a locally common dynamical
mode.  Under that contract, delay becomes phase, the weighted snapshot becomes
a phasor sum, and an exact coherence audit follows.

\section{Preregistration and scope}

The hypotheses, assumptions, falsifiers, and verification program were fixed
before the manuscript was drafted.

\subsection{Hypotheses}

For a locally common harmonic mode and known channel delays:

\begin{enumerate}
\item the aligned signal fraction is the squared magnitude of the weighted
delay phasor;
\item its loss has an exact pairwise cosine decomposition;
\item the loss is bounded by frequency squared times weighted delay variance;
\item this bound is second-order sharp as delay spread tends to zero; and
\item a bandlimited common-mode process inherits a spectral version of the
same guarantee.
\end{enumerate}

\subsection{Assumptions}

The channels share a locally common mode after centering and scaling.  Effective
weights are nonnegative and normalized, delays are known relative to a declared
reference time, and the local response is linear.  Measurement noise, spatial
mismatch, missingness, taxonomic bias, calibration drift, and channel-specific
ecological dynamics are held outside the timing term.  These exclusions are
substantive: a real deployment must test them rather than inherit them.

\subsection{Falsifiers and controls}

The computational gate would fail if the pairwise identity exceeded numerical
tolerance, if $\eta$ left $[0,1]$, if the variance bound were violated, or if
independent time-domain estimates failed to converge to the phasor prediction.
Negative controls included equal delays, a static mode, one active channel,
balanced half-period cancellation, and full-period phase wrapping.

\section{Model}

Let the centered common ecological perturbation be
\begin{equation}
x(t)=\Re\{z e^{\mathrm{i}\omega t}\},
\label{eq:mode}
\end{equation}
where $z$ is a complex amplitude and $\omega$ is angular frequency.  Channel
$i$ observes the same local mode at delay $\tau_i$.  Channel sensitivity can be
absorbed into an effective weight $w_i$ after sign and scale have been validated.
We assume
\begin{equation}
w_i\ge0, \qquad \sum_{i=1}^n w_i=1.
\label{eq:weights}
\end{equation}
The fused response is
\begin{equation}
y(t)=\sum_i w_i x(t-\tau_i)
=\Re\{z e^{\mathrm{i}\omega t}A(\omega)\},
\end{equation}
with transfer phasor
\begin{equation}
A(\omega)=\sum_i w_i e^{-\mathrm{i}\omega\tau_i}.
\label{eq:phasor}
\end{equation}
We define snapshot alignment efficiency as the common-mode output-to-input
energy ratio
\begin{equation}
\eta(\omega)=|A(\omega)|^2.
\label{eq:eta}
\end{equation}
This quantity measures temporal coherence only.  It is not an information
content, accuracy, causality, or conservation-value score.

Define the weighted mean delay and variance by
\begin{equation}
\bar\tau=\sum_i w_i\tau_i, \qquad
\Var_w(\tau)=\sum_i w_i(\tau_i-\bar\tau)^2.
\label{eq:variance}
\end{equation}
Only relative timing should matter: translating every $\tau_i$ by the same
constant rotates $A$ without changing $\eta$.

\section{Results}

The formal artifact \texttt{SnapshotAlignment.lean} defines the weighted
complex \texttt{phasor}, its squared-norm \texttt{alignment}, the
\texttt{weightedMean}, and \texttt{delayVariance}.  It contains exactly six
audited theorems: \path{phasor_energy_pairwise_identity},
\path{alignment_efficiency_unit_interval},
\path{delay_variance_bounds_alignment_loss},
\path{alignment_budget_sufficient_for_band},
\path{equal_delay_perfect_alignment}, and
\path{two_channel_half_period_cancellation}.  The first is stated over
an arbitrary finite index type without weight hypotheses.  The next three use
nonnegative normalized weights; the band theorem additionally assumes
$0<\Omega$, $0\le q\le1$, $|\omega|\le\Omega$, and the displayed variance
budget.  The equal-delay theorem assumes a nonempty finite channel type and
weights summing to one.  The cancellation theorem assumes exactly two equal
weights and the explicit antipodal-phase equation.  These are mathematical
claims about the defined model only.

\subsection{Exact pairwise identity}

\textbf{Proposition 1 (snapshot alignment identity).}
Under Eqs.~\eqref{eq:weights}--\eqref{eq:eta},
\begin{align}
\eta(\omega)
&=\sum_i w_i^2+2\sum_{i<j}w_iw_j
\cos\!\left(\omega(\tau_i-\tau_j)\right),
\label{eq:pair}\\
1-\eta(\omega)
&=2\sum_{i<j}w_iw_j
\left[1-\cos\!\left(\omega(\tau_i-\tau_j)\right)\right].
\label{eq:loss}
\end{align}

\emph{Derivation.}  Expand $A\overline{A}$ into diagonal and off-diagonal
terms.  Pair each complex exponential with its conjugate to obtain a cosine.
Then subtract the result from
$1=(\sum_iw_i)^2=\sum_iw_i^2+2\sum_{i<j}w_iw_j$.

Equation~\eqref{eq:loss} assigns every unit of modeled coherence loss to a
pair of channels.  It also makes the bounds $0\le\eta\le1$ immediate: squared
magnitude gives the lower bound and each term in the loss is nonnegative.

\subsection{Delay-variance bound}

\textbf{Theorem 1 (delay variance bounds alignment loss).}
For every real $\omega$,
\begin{equation}
1-\eta(\omega)\le\omega^2\Var_w(\tau),
\qquad
\eta(\omega)\ge1-\omega^2\Var_w(\tau).
\label{eq:mainbound}
\end{equation}

\emph{Derivation.}  Apply $1-\cos u\le u^2/2$ termwise in
Eq.~\eqref{eq:loss}, then use the weighted pairwise-variance identity
\begin{equation}
\sum_{i<j}w_iw_j(\tau_i-\tau_j)^2=\Var_w(\tau).
\end{equation}
For small delays, $1-\cos u=u^2/2+O(u^4)$, so the ratio
$(1-\eta)/(\omega^2\Var_w)$ tends to one whenever the delay shape has nonzero
variance.  The gate recovered second-order convergence under repeated halving.

\subsection{Bandlimited process and operational rule}

The following spectral extension is an ordinary analytic consequence and was
not formalized among the six Lean theorems.  Suppose a wide-sense stationary
common-mode process has spectral density
$S_x(\omega)$ supported on $|\omega|\le\Omega$.  By linear filtering and
Parseval's relation, its normalized fused common-mode energy is the
$S_x$-weighted average of $\eta(\omega)$.  Equation~\eqref{eq:mainbound}
therefore yields
\begin{equation}
\frac{\int S_x(\omega)\eta(\omega)\,d\omega}
{\int S_x(\omega)\,d\omega}
\ge 1-\Omega^2\Var_w(\tau).
\label{eq:spectral}
\end{equation}

\textbf{Corollary 1 (alignment budget).}  If a monitoring product declares a
maximum modeled frequency $\Omega$ and minimum acceptable alignment $q\in[0,1]$,
then
\begin{equation}
\boxed{\Var_w(\tau)\le\frac{1-q}{\Omega^2}}
\label{eq:budget}
\end{equation}
is sufficient for $\eta(\omega)\ge q$ throughout the declared band.

This rule is conservative.  It can be audited from timestamps and weights, but
it is only meaningful if $\Omega$ is scientifically justified.  Declaring a
slower process makes the budget looser; that choice must be tied to an external
ecological hypothesis rather than selected to pass the gate.

\subsection{Controls and edge cases}

\begin{center}
\small
\textbf{Table 1. Analytic controls defining the scope of the coherence claim.}\\[2pt]
\begin{tabular}{p{0.29\linewidth}p{0.24\linewidth}p{0.37\linewidth}}
\toprule
Case & Result & Interpretation \\
\midrule
All $\tau_i$ equal & $\eta=1$ & Common timestamp preserves the mode. \\
$\omega=0$ & $\eta=1$ & Static differences are unaffected by age. \\
One active channel & $\eta=1$ & Coherence alone does not imply completeness. \\
Two equal channels, half-cycle apart & $\eta=0$ & An average can erase the common signal. \\
Two equal channels, full-cycle apart & $\eta=1$ & Raw delay variance is sufficient, not necessary. \\
\bottomrule
\end{tabular}
\end{center}

For two channels with weights $w$ and $1-w$ and delay difference $\Delta\tau$,
Eq.~\eqref{eq:loss} reduces to
\begin{equation}
\eta=1-4w(1-w)\sin^2\!\left(\frac{\omega\Delta\tau}{2}\right).
\end{equation}
The worst cancellation occurs at balanced weights and an odd half-cycle.  In
contrast, an integer full-cycle difference is phase aligned.  This periodicity
is why Eq.~\eqref{eq:budget} must never be represented as a necessary condition.

\section{Verification}

The exact program \texttt{verification.py} uses only the Python standard
library and a fixed seed.  It was written and executed before this manuscript.
The captured, untruncated output is \texttt{verification\_output.txt}.

The gate passed 21 of 21 checks.  It included a hand-auditable three-channel
case, 5,000 randomized cases with two to twelve channels, weights spanning seven
orders of magnitude, delays on both sides of the reference time, and frequencies
spanning more than four orders of magnitude.  The maximum pairwise-identity
error was $3.531\times10^{-14}$.  The largest computed excess above the variance
bound was $1.178\times10^{-16}$, consistent with floating-point roundoff.

An independent time-domain calculation sampled delayed cosines at 256, 1,024,
4,096, and 16,384 points and reproduced the phasor energy to below
$2.5\times10^{-15}$.  The small-delay experiment observed convergence orders
1.9985, 1.9996, and 1.9999.  A four-frequency mixture passed the spectral bound,
and a grid of 1,001 frequencies passed the operational $q=0.90$ budget.

\begingroup\raggedright
These checks reduce implementation and algebra risk but are not the formal
proof.  Formal support comes from the separate audited Lean artifact.  The
numeric suite does not validate the common-mode assumption, establish novelty,
or demonstrate field utility.  The small-delay convergence and spectral-mixture
experiments are numeric checks of claims not separately stated as Lean
theorems.\par
\endgroup

\section{Empirical verification plan}

A field test should begin with a named common-mode hypothesis, not with the
composite itself.  For example, a monitoring program might preregister that
several normalized indicators respond with positive sensitivity to a seasonal
productivity mode over a specified frequency band.  Training years would be
used to estimate signs, gains, delay distributions, and $\Omega$.  Weights,
$q$, and interpolation rules would then be frozen.

On held-out years, the program would compare: (i) the original asynchronous
composite; (ii) a common-time reconstruction using only training-fitted
state-space or interpolation models; and (iii) a simple maximum-age baseline.
The central test is whether snapshots inside the variance budget achieve
coherence at least $q$ against an independent common-mode reference, with delay
uncertainty propagated.  Snapshots outside the budget must remain in the test
set as predicted failures.

The model should be rejected for the intended application if the channels lack
the preregistered common mode, if estimated signs differ, if alignment does not
improve held-out coherence, if the bound fails after delay uncertainty is
included, or if the simpler baseline predicts failures equally well.  Genuine
ecological phase differences should be modeled as biology, not ``corrected''
away to make a score smoother.

\section{Discussion}

The main contribution is not the observation that timestamps matter.  It is a
compact contract connecting three quantities a monitoring product can declare:
the frequency band it claims to represent, the coherence it promises, and the
weighted dispersion of the timestamps it actually uses.  The pairwise identity
also supports diagnosis: channels with large $w_iw_j$ and large phase separation
dominate predicted loss, so remediation can target the specific stale pair
rather than refresh every stream.

Run 118 is distinct from the prior D-Score lineage.  Measurement closure asks
whether a value is statistically resolvable.  Monitoring water-filling asks
where finite observation power should go.  Revisit scheduling asks when each
site should be observed again.  Manipulation-proof weighting asks how strategic
elasticity should affect weights.  Snapshot alignment asks whether the chosen
values, weights, and ages can represent one common-time state.  A mature system
would need all of these gates; passing one does not imply another.

The result also supplies a useful negative lesson.  ``Fresh enough'' cannot be
defined by a universal maximum age.  A one-month mismatch may be negligible for
a decadal mode and destructive for a seasonal or disturbance mode.  Conversely,
large age differences can phase-wrap for a single exact sinusoid, although such
coincidence is fragile under broadband dynamics.  The declared ecological
bandwidth, not the calendar alone, is the essential comparison scale.

\section{Limitations}

The analysis is intentionally narrow.  A composite may be temporally coherent
yet biased, gamed, underpowered, spatially mismatched, or unrelated to ecological
value.  Nonnegative weights exclude validated anticorrelated channels.  A single
common mode excludes component-specific dynamics.  Harmonic analysis is local
and does not by itself handle regime shifts, hysteresis, irregular missingness,
aliasing, nonlinear response, or uncertain clocks.  The bound becomes negative
and uninformative when $\Omega^2\Var_w(\tau)>1$.

Most importantly, real ecological asynchrony is not automatically a defect.
Phenological divergence can encode interactions, adaptation, or stability.  The
proposed gate applies only when channels are contractually intended to track the
same local mode.  No external dataset was analyzed, no field performance is
claimed, and the targeted literature search cannot establish novelty or
priority.

\section{Reproducibility and declarations}

\textbf{Reproducibility.}  Run \texttt{python3 verification.py} with Python
3.11 or later; no third-party package or external data is required.  The seed,
tolerance, hypotheses, assumptions, and falsifiers are printed in the captured
output.  The manuscript compiles from \texttt{paper.tex} with Tectonic 0.17.0.

\textbf{Data and code availability.}  All code and synthetic parameters used
in this generated exploration are contained in the immutable Run 118 package.
No external ecological data were generated or analyzed.

\textbf{Funding.}  No specific external funding supported this run.

\textbf{Conflicts of interest.}  Justin D. Hart is founder of Viridis LLC,
which may develop biodiversity measurement and assurance products.  The
separate business note reports zero signed revenue and no named customer
validation attributable to this result.

\textbf{Accountable authorship.}  Justin D. Hart is the accountable human
author and retains responsibility for scientific claims, corrections, and any
future submission.  No AI system is represented as a journal author.

\textbf{Substantive AI-use disclosure.}  Under Justin D. Hart's direction,
OpenAI Codex searched the local Viridis corpus and current literature, proposed
and preregistered the hypothesis, wrote and executed the verification code,
drafted the manuscript, assembled the package, and reconciled this final version
line by line against the landed formal artifact.  Historical Claude-authored
Viridis artifacts were read for context with attribution preserved.  Aristotle
served as the automated theorem-proving system for the six frozen Lean targets;
Codex then independently rebuilt and audited the returned artifact for statement
freeze, proof holes, forbidden constructs, allowed axioms, non-vacuity, and
significance.  AI systems are tools and contributors, not authors or creators.
No external peer review has occurred.

\textbf{Stage and route.}  This version is the exact
\texttt{FINAL\_RECONCILED} manuscript bound to the local package.
The authoritative ledger row is \texttt{null}: no exact Run 118 identity occurs
in the root ledger.  The advisory route is working corpus.  The result is a
domain-applying biodiversity-measurement theorem and is not admitted to the
spine (\texttt{spine\_admitted=false}).  These local states confer no ledger
approval, publication authorization, peer acceptance, adoption, or revenue.

\begin{thebibliography}{9}
\bibitem{piatkowski2025}
J. Piatkowski, L. Karbowiak, and F. Depta,
``A lightweight algorithm for synchronized multimodal data acquisition using
temporal sample alignment,'' \emph{Scientific Reports} \textbf{15}, 22711
(2025). \href{https://doi.org/10.1038/s41598-025-07797-7}{doi:10.1038/s41598-025-07797-7}.

\bibitem{terasakihart2025}
D. E. Terasaki Hart, T. N. Bui, L. Di Maggio, et al.,
``Global phenology maps reveal the drivers and effects of seasonal asynchrony,''
\emph{Nature} \textbf{645}, 133--140 (2025).
\href{https://doi.org/10.1038/s41586-025-09410-3}{doi:10.1038/s41586-025-09410-3}.

\bibitem{lang2025}
W. Lang, Y. Zhang, X. Li, et al.,
``Phenological divergence between plants and animals under climate change,''
\emph{Nature Ecology \& Evolution} \textbf{9}, 261--272 (2025).
\href{https://doi.org/10.1038/s41559-024-02597-0}{doi:10.1038/s41559-024-02597-0}.

\bibitem{leung2024}
B. Leung and A. Gonzalez,
``Global monitoring for biodiversity: Uncertainty, risk, and power analyses to
support trend change detection,'' \emph{Science Advances} \textbf{10},
eadj1448 (2024).
\href{https://doi.org/10.1126/sciadv.adj1448}{doi:10.1126/sciadv.adj1448}.

\bibitem{schulte2018}
H. Schulte to Buehne and N. Pettorelli,
``Better together: Integrating and fusing multispectral and radar satellite
imagery to inform biodiversity monitoring, ecological research and
conservation science,'' \emph{Methods in Ecology and Evolution} \textbf{9},
849--865 (2018).
\href{https://doi.org/10.1111/2041-210X.12942}{doi:10.1111/2041-210X.12942}.
\end{thebibliography}

\end{document}
